% Supplement to A.4.42.  Root may merge deliberately; no existing source edited.
\paragraph{The discrepancy realizes the entire signed-period kernel.}
Retain the actual maps and markings (FCD.1)--(FCD.2), and let
\[
 \varpi:J_c\longrightarrow\operatorname{SL}_2(\mathbb Z/12)
 \tag{CDF.1}
\]
be the signed ordinary compact-period map of A.4.22.
It is onto, and its kernel has order $3^4$: the actual
ordinary image has order $1152\cdot3^4$, while the
period image has order $1152$.
Define the subgroup
\[
 \mathcal K^\dagger=
 C_{G^\dagger}(H^\dagger)\cap
 \operatorname{im}(\mathcal M_0)\cap\ker j_0^\dagger.
 \tag{CDF.2}
\]
Then the exact image, with the same $d$ and projection as
in (FCD.4), is
\[
 d\bigl(\operatorname{pr}_G(\mathcal K^\dagger)\bigr)
       =\ker\varpi\cong C_3^4.
 \tag{CDF.3}
\]
This equality remains valid after adjoining any of the
finite solvable families allowed in A.4.42 and imposing
identity on those extra coordinates.

\begin{proof}
We first give actual lifts of all six first-order
directions, rather than use their dimension as a lifting
argument. Put
\[
 A=\mathbb F_3[H]
   =\mathbb F_3[X,Y]/(X^3,Y^3),\qquad
 X=h_1-1,\quad Y=h_2-1,\quad\mathfrak m=(X,Y).
 \tag{CDF.4}
\]
The actual module $M$ is free of rank two over $A$ by
A.4.20. The involution is $h\mapsto h^{-1}$, so
$X^*=-X$ and $Y^*=-Y$ modulo $\mathfrak m^2$.
Use its actual nondegenerate Hermitian form and denote
the Hermitian adjoint of an $A$-linear endomorphism by
$\dagger$. In the retained free basis its augmentation
form is $H_0=2I_2$, as independently checked in A.4.21.

Choose any two matrices $D,E\in\operatorname{Mat}_2(\mathbb F_3)$
satisfying
\[
 D^tH_0=H_0D,\qquad E^tH_0=H_0E.
 \tag{CDF.5}
\]
Each is symmetric in these coordinates and has three
independent coefficients. Extend them as constant
$A$-linear matrices in the chosen free $A$-basis and set
\[
 W=2(XD+YE),\qquad
 Z=\frac{W-W^\dagger}{2},\qquad
 g_{D,E}=(I+Z)(I-Z)^{-1}.
 \tag{CDF.6}
\]
Here all matrices are endomorphisms of the actual
eighteen-dimensional module. Because
$\mathfrak m^5=0$, the inverse in (CDF.6) is the
finite sum $I+Z+Z^2+Z^3+Z^4$.
One has $Z^\dagger=-Z$, and consequently
$g_{D,E}^\dagger g_{D,E}=I$ by direct multiplication
of these polynomials in $Z$.
The construction is $A$-linear, so $g_{D,E}$ commutes
with the two named deck generators.

For the exact reduction, (CDF.5) and the displayed
involution give $W^\dagger\equiv-W\pmod{\mathfrak m^2}$.
Thus $Z\equiv W$ and
\[
 \overline g_{D,E}
       =I+2W=I+XD+YE
             \quad\text{on }M/\mathfrak m^2M.
 \tag{CDF.7}
\]
The full $g_{D,E}$ belongs to
$1+\operatorname{Mat}_2(\mathfrak m)$.
Its ninth power is identity, since
$(g_{D,E}-I)^5=0$ in characteristic three.
It therefore has odd order; both its determinant
and its spinor norm are one. Hence every lift in
(CDF.6) is in $C_{\Omega^-_{18}(3)}(H)$.
This proves actual liftability of all six parameters.

The reductions (CDF.7) form an elementary abelian
group $\mathcal U\cong C_3^6$, because the product
of two first-order terms vanishes modulo
$\mathfrak m^2$. Pair each reduction with the
identity signed period. These pairs belong to the
exact normalizer group (NC.1)--(NC.2): their label
map is identity and their spinor is one.
Thus A.4.22 supplies their actual map to $J_c$.
Its inner kernel is the stated eighteen-element
group $\{\pm h:h\in H\}$. Within $\mathcal U$,
the constant coefficient is $I_2$, so the negative
coset is absent. The remaining nine elements are
exactly
\[
 I+aXI_2+bYI_2,\qquad a,b\in\mathbb F_3.
 \tag{CDF.8}
\]
The image of $\mathcal U$ in $J_c$ is consequently
$C_3^{6-2}=C_3^4$ and has identity period.
The order in (CDF.1) now identifies that image
with the entire $\ker\varpi$.
In particular this argument computes the group
structure of the period kernel, in addition to
its order.

For each actual $g_{D,E}$, the perfect-image argument
(FCD.8)--(FCD.9), with the same lower and all nine
forgetting coordinates, supplies an actual pure word
whose multiplicity matrix is $g_{D,E}$ and whose other
coordinates are identity. Its image centralizes
$H^\dagger$ by the compatible-pair reconstruction,
and belongs to $\mathcal K^\dagger$.
Formula (TF.15a) maps it to the class of (CDF.7).
Thus the left side of (CDF.3) contains
$\ker\varpi$.

For the reverse inclusion retain the literal compact
period map of the source. Write $C$ for the period
of an actual representative in the primitive upstairs
basis $(b^3,a^3)$. Forgetting marked points preserves
this homology map, so $j_0^\dagger=1$ implies
$C=I\pmod{12}$.
The period used by $d$ in (NC.1)--(NC.2) is
$q_{\rm down}=JCJ$, with
$J=\left(\begin{smallmatrix}0&1\\1&0\end{smallmatrix}\right)$,
as proved in (NJ.5)--(NJ.7) and retained in A.4.25.
It too is identity. Therefore every element in the
left side of (CDF.3) has trivial period, proving
the reverse inclusion without assuming equality of
$d$ and $j_0^\dagger$ themselves.
The same proof works with the additional solvable
coordinates: join them to (FCD.8) and increase the
derived depth exactly as in A.4.42.

The independent program
\path{work/ghost_code/check_cm_discrepancy_family.py}
constructs an actual free $A$-basis lifting the
retained two constant quotient generators. It
computes all six lifts (CDF.6), their full
eighteen-dimensional matrices, their orthogonality,
deck commutation, and ninth powers, and verifies
the exact quotient intertwining for each one.
Their six coefficient vectors span the same space
as the six actual geometric directions retained
in A.4.21. The program computes the two actual
inner translation rows, of rank two, and the
rank-four quotient map
\[
 (D,E)\longmapsto
 (D_{11}-D_{22},D_{12},E_{11}-E_{22},E_{12}).
 \tag{CDF.9}
\]
It enumerates all $729$ first-order values and
their $81$ outer classes, each with exactly nine
representatives. All matrices, ranks, input hashes,
and the actual execution output are retained.
\end{proof}
